\chapter[Computer Simulation of Time-Correlated Processes]{Computer Simulation of Time-Correlated Random Processes}
\label{appendix:time_correlated_processes}

When generating channel realizations for a time-varying channel, it is necessary to simulate a zero-mean complex Gaussian random process with a prescribed correlation in time.
This appendix presents two practical methods for simulating time-correlated random processes on a computer.
The first method uses the autocorrelation function as input and is based on the Cholesky decomposition of the covariance matrix.
The second method uses the power spectral density as input and is based on filtering a white Gaussian process with a linear time-invariant system.
In the simulation framework developed for this dissertation, both methods are employed to generate time-correlated channel realizations.


\section{Cholesky Decomposition Method}
\label{appendix:cholesky_method}
% =====================================
%   Cholesky Decomposition Method
% =====================================

The Cholesky decomposition method for generating a zero-mean unit-variance complex Gaussian process $\mathbf{h}$ of length $N$ with a specified autocorrelation function $R_h(\tau)$ consists of the following steps:
\begin{itemize}[label={}, leftmargin=0.5cm]
    
    \item \textbf{Step 1:} 
    Build the $N \times N$ covariance matrix $\mathbf{C}$ with entries $C_{i,j} = R_h((i-j)T_s)$.
    
    \item \textbf{Step 2:} 
    Compute the Cholesky decomposition of the covariance matrix $\mathbf{C} = \mathbf{L}\mathbf{L}^H$, where $\mathbf{L}$ is a lower triangular matrix.
    
    \item \textbf{Step 3:} 
    Generate a column vector $\mathbf{w}$ of $N$ i.i.d. white complex Gaussian random variables with zero mean and unit variance.
    
    \item \textbf{Step 4:}
    Obtain the desired Gaussian process as $\mathbf{h} = \mathbf{L}\mathbf{w}$.

\end{itemize}

We show that the Cholesky decomposition method effectively generates a zero-mean unit-variance complex Gaussian process $\mathbf{h}$ of length $N$ with a specified autocorrelation function $R_h(\tau)$.

First, let us verify that the generated process $\mathbf{h}$ is a complex Gaussian process with zero mean and the specified autocorrelation function.

Each element $h_n = \sum_{k=1}^n L_{k, n} w_k$ of the process $\mathbf{h}$ is a complex Gaussian random variable, since $w_k$ are \ac{i.i.d.} complex Gaussian random variables and the linear combination of complex Gaussian random variables is itself also a complex Gaussian random variable. Let us now calculate the mean and covariance matrix of the generated process $\mathbf{h}$.

The expected value of each element $h_n$ is given by:
\begin{equation}
\label{eq:cholensky_mean_h}
    \mathbb{E}[h_n] = \mathbb{E}\left[\sum_{k=1}^n L_{k, n} w_k\right] = \sum_{k=1}^n L_{k, n} \mathbb{E}[w_k] = 0
\end{equation}

The covariance matrix of $\mathbf{h}$ is given by:
\begin{equation}
\label{eq:cholensky_covariance_h}
\begin{aligned}
    \text{Cov}[\mathbf{h}] &= \mathbb{E}[(\mathbf{h} - \mathbb{E}[\mathbf{h}])(\mathbf{h} - \mathbb{E}[\mathbf{h}])^H] = \mathbb{E}[\mathbf{h}\mathbf{h}^H] \\
    &= \mathbb{E}[(\mathbf{L}\mathbf{w})(\mathbf{L}\mathbf{w})^H] = \mathbb{E}[\mathbf{L}\mathbf{w}\mathbf{w}^H\mathbf{L}^H] = \mathbf{L} \,  \mathbb{E}[\mathbf{w}\mathbf{w}^H] \, \mathbf{L}^H \\
    &= \mathbf{L} \, \mathbf{I}_N \, \mathbf{L}^H = \mathbf{L} \mathbf{L}^H = \mathbf{C}
\end{aligned}
\end{equation}

From~\eqref{eq:cholensky_mean_h} and~\eqref{eq:cholensky_covariance_h}, the mean of the generated process $\mathbf{h}$ is zero and the covariance of $\mathbf{h}$ matches the specified autocorrelation function $R_h(\tau)$ by construction. This confirms that the Cholesky decomposition method effectively generates a zero-mean unit-variance complex Gaussian process with the desired correlation properties:
\begin{equation}
    \mathbf{h} \sim \mathcal{CN}(\mathbf{0}, \mathbf{C})
\end{equation}

Second, let us demonstrate that the Cholesky decomposition of the covariance matrix $\mathbf{C}$ always exists.
It is well-known that the Cholesky decomposition of a matrix $\mathbf{C}$ exists and is unique if and only if $\mathbf{C}$ is a Hermitian matrix and a positive-definite matrix, i.e.:
\begin{equation}
    \exists ! \, \mathbf{L} : \mathbf{C} = \mathbf{L}\mathbf{L}^H \iff \mathbf{C} = \mathbf{C}^H \quad \text{and} \quad \forall \mathbf{x} \in \mathbb{C}^N \setminus \{\mathbf{0}\} : \, \mathbf{x}^H \mathbf{C} \mathbf{x} > 0
\end{equation}

Consider an autocorrelation function $R_h(\tau)$ of a wide-sense stationary process, i.e.\ $R_h$ does not depend on the absolute time $t$, but only on the time difference $\tau$.
Therefore, the covariance matrix $\mathbf{C}$ is a real and symmetric Toeplitz matrix by construction. Thus, $\mathbf{C}$ is a Hermitian matrix.

We prove that $\mathbf{C}$ is positive-definite by contradiction.
First, let us simplify the expression $\mathbf{x}^H \mathbf{C} \mathbf{x}$:
\begin{align}
    \mathbf{x}^H \mathbf{C} \mathbf{x} &= \sum_{i=1}^N \sum_{j=1}^N x_i^* C_{i,j} x_j = \sum_{i=1}^N \sum_{j=1}^N x_i^* R_h((i-j)T_s) x_j \notag \\
    &= \sum_{i=1}^N \sum_{j=1}^N x_i^* \left( \int_{-\infty}^{\infty} S_h(f) e^{j 2 \pi f (i-j)T_s} df \right) x_j \notag \\
    &= \int_{-\infty}^{\infty} S_h(f) \left( \sum_{i=1}^N x_i^* e^{j 2 \pi f i T_s} \right) \left( \sum_{j=1}^N x_j e^{-j 2 \pi f j T_s} \right) df \displaybreak \notag \\
    &= \int_{-\infty}^{\infty} S_h(f) \, \left|\sum_{i=1}^N x_i e^{-j 2 \pi f i T_s}\right|^2 \; df \notag \\
    &= \int_{-\infty}^{\infty} S_h(f) \, |X(f)|^2 \; df
    \label{eq:cholensky_positive_definite_C}
\end{align}

We consider an autocorrelation function $R_h(\tau)$ that corresponds to a power spectral density that makes physical sense, i.e.\ $S_h(f)$ is positive for all frequencies $f$ and strictly positive for a set of positive measure.
Then, it is easy to see that $\mathbf{C}$ is positive semi-definite, since the integrand is non-negative for all frequencies $f$.

Assume that $\mathbf{C}$ is not positive-definite. Then, there exists a non-zero vector $\mathbf{x} \in \mathbb{C}^N$ such that $\mathbf{x}^H \mathbf{C} \mathbf{x} = 0$.
For that to be true, the integrand must be zero, which implies that $X(f) = 0$ for all frequencies $f$ where $S_h(f) > 0$, which is a set of positive measure. 
$X(f)$ can be seen as the $\mathcal{Z}$-transform of $\mathbf{x}$ evaluated on the unit circle, which corresponds to a complex  polynomial of degree at most $N-1$. 
By the Fundamental Theorem of Algebra, such a nonzero polynomial has at most $N-1$ roots. Therefore, $\mathbf{x}$ must be the zero vector, which contradicts the assumption. 
Hence, $\mathbf{C}$ is positive-definite.

\section{Linear Filtering Method}
\label{appendix:psd_method}
% ===============================
%   Linear Filtering Method
% ===============================

The linear filtering method for generating a zero-mean unit-variance complex Gaussian process with a specified power spectral density $S_h(f)$ consists of the following steps:
\begin{itemize}[label={}, leftmargin=0.5cm]

    \item \textbf{Step 1:}
    Evaluate the target \ac{PSD} $S_h(f)$ on an \acs{FFT} frequency grid of $N_g$ uniformly spaced normalized frequencies $\nu_k \in [-1/2, 1/2)$.
    To avoid aliasing, the sampling period $T_s$ must satisfy the Nyquist condition $f_D T_s < 1/2$, where $f_D$ denotes the maximum Doppler frequency.

    \item \textbf{Step 2:}
    Construct the frequency-domain shaping response by taking the nonnegative square root of the sampled \ac{PSD}: $G[\nu_k] = \sqrt{S_h[\nu_k]}$.

    \item \textbf{Step 3:}
    Apply an inverse discrete Fourier transform to the shaping response $G[\nu_k]$ in order to obtain the impulse response of a finite-length \acs{FIR} filter, $g[n] = \mathcal{F}^{-1}\{G[\nu_k]\}$, and normalize this filter to unit energy.

    \item \textbf{Step 4:}
    Generate a sequence of white proper complex Gaussian random variables $w[n]$ with zero mean and unit variance per sample.

    \item \textbf{Step 5:}
    Obtain the desired time-correlated output process by convolving the white-noise sequence with the shaping filter: $h[n] = \sum_m g[m]\, w[n-m]$.

\end{itemize}

The underlying principle behind this method is that passing white complex Gaussian noise through a linear time-invariant filter with frequency response $G(f)$ yields an output process whose power spectral density is equal to $|G(f)|^2$.
Consequently, by choosing $G(f) = \sqrt{S_h(f)}$, the output process is shaped so as to have the desired power spectral density.
The derivation below makes this argument precise.

First, note that the output process $h(t)$ is a linear transformation of the white Gaussian input process $w(t)$ and is therefore itself a complex Gaussian process.
Moreover, because the input process has zero mean and the filter is linear, the mean of the output process is also zero.

Let us now determine the autocorrelation function of the output process $h(t)$:
\begin{fleqn}[0.75cm]
\begin{align}
    R_h(\tau) 
    &= \mathbb{E}[h(t) \, h^*(t-\tau)] \notag \\
    &= \mathbb{E}\left[ \int_{-\infty}^{+\infty} g(u) \, w(t-u) \, du \, \int_{-\infty}^{+\infty} g(v) \, w^*(t-\tau-v) \, dv \right] \notag \\
    &= \int_{-\infty}^{+\infty} \int_{-\infty}^{+\infty} g(u) \, g^*(v) \, \mathbb{E}\big[w(t-u) \, w^*(t-\tau-v)\big] \, du \, dv \notag \\
    & \qquad \bullet \quad \mathbb{E}\big[w(t-u) \, w^*(t-\tau-v)\big] \notag \\
    & \qquad \quad \quad = \mathbb{E}\big[w(t-u) \, w^*((t-u) - (\tau-v-u))\big] = R_w(\tau + v - u) = \delta(\tau + v - u) \notag \\
    &= \int_{-\infty}^{+\infty} \int_{-\infty}^{+\infty} g(u) \, g^*(v) \, \delta(\tau + v - u) \, du \, dv \notag \\
    &= \int_{-\infty}^{+\infty} g(u) \, g^*(u-\tau) \, du \notag \\
    &= R_g(\tau) \label{eq:psd_autocorrelation_output_h}
\end{align}
\end{fleqn}

Equation~\eqref{eq:psd_autocorrelation_output_h} shows that the autocorrelation function of the output process is equal to the autocorrelation function of the shaping filter.
In particular, because the filter is normalized to unit energy, the variance of the output process is given by $R_h(0) = R_g(0) = \int_{-\infty}^{+\infty} |g(u)|^2 du = 1$.
Hence, the generated process has unit variance.

Next, we compute the Fourier transform of the autocorrelation function in order to obtain the power spectral density of the output process:
\begin{fleqn}[0.75cm]
\begin{align}
    S_h(f) 
    &= \int_{-\infty}^{+\infty} R_g(\tau) \, e^{-j 2 \pi f \tau} \, d\tau \notag \\
    &= \int_{-\infty}^{+\infty} \left( \int_{-\infty}^{+\infty} g(u) \, g^*(u-\tau) \, du \right) \, e^{-j 2 \pi f \tau} \, d\tau \notag \\
    &= \int_{-\infty}^{+\infty} \int_{-\infty}^{+\infty} g(u) \, g^*(u-\tau) \, e^{-j 2 \pi f \tau} \, du \, d\tau \notag \\
    &= \int_{-\infty}^{+\infty} \int_{-\infty}^{+\infty} g(u) \, g^*(v) \, e^{-j 2 \pi f (u-v)} \, du \, dv \notag \\
    &= \left( \int_{-\infty}^{+\infty} g(u) \, e^{-j 2 \pi f u} \, du \right) \, \left( \int_{-\infty}^{+\infty} g^*(v) \, e^{j 2 \pi f v} \, dv \right) \notag \\
    &= G(f) \, G^*(f) \notag \\
    &= \left|G(f)\right|^2 \label{eq:psd_output_h}
\end{align}
\end{fleqn}

Equation~\eqref{eq:psd_output_h} confirms that the power spectral density of the output process equals the squared magnitude of the filter's frequency response.
Therefore, choosing the shaping filter such that $G(f) = \sqrt{S_h(f)}$ ensures that the generated process has the desired power spectral density.

In a practical computer implementation, the process cannot be generated in continuous time and with an infinite-length filter.
Instead, the shaping filter is approximated by a finite-length discrete-time filter obtained from sampled values of the target \ac{PSD} via the inverse \ac{DFT}.
It is important to choose the filter length $N_g$ sufficiently large, so that truncation effects become negligible and the finite-length filter accurately approximates the target correlation properties of the process.

